% GENERATED FILE. Edit canonical lessons, then run npm run generate:books.
% canonical-source-hash: fnv1a64:82044303a624e39d
% canonical-lessons: JA-C510-guuzen, JA-C510-kikkake, JA-C510-futo, JA-C510-ikasu, JA-C510-megumareru

\chapter{By Chance}
\label{ch:ja-by-chance}
\hlchaptermodality{510}

\begin{chapteropening}
\textbf{By the end of this chapter:} \emph{I can say what I would do if a chance came.}

Complete the last lesson of chapter 510: I can say what I would do if a chance came.
\end{chapteropening}

\section[\texorpdfstring{\ja{ぐうぜん} (gūzen)}{ぐうぜん (gūzen)}]{\ja{ぐうぜん} (gūzen) — by chance; a coincidence}

\label{lesson:JA-C510-guuzen}



\begin{quote}
Before the new one: say the Japanese for probability, then the Japanese for to hit (the mark).
\end{quote}

\subsection*{You'll want to know: \ja{ぐうぜん}}
\textbf{\ja{ぐうぜん}} — \emph{gūzen} — \textquotedblleft{}by chance; a coincidence\textquotedblright{}.

\begin{quote}
\textbf{\ja{もし} \ja{ぐうぜん} \ja{かれに} \ja{あったら、てがみを} \ja{わたして} \ja{ください。}} — \emph{moshi gūzen kare ni attara, tegami o watashite kudasai} — \textquotedblleft{}If you happen to meet him, please give him the letter.\textquotedblright{}
\end{quote}

\begin{cousinweb}
\textbf{\ja{ぐうぜん}} (\emph{gūzen}) is Sino-Japanese (kango): two Chinese characters read with their on'yomi, readings borrowed from Chinese. They are \emph{gū}, \textquotedblleft{}accidentally\textquotedblright{} (Middle Chinese \emph{nguwH}) and \emph{zen}, \textquotedblleft{}sort of thing\textquotedblright{} (\emph{nyen}) (KANJIDIC2; Middle Chinese in Baxter's transcription).
\end{cousinweb}

\begin{grammarlens}[title={What you've built}]
One more word for saying what happens if.
\end{grammarlens}

\subsection*{Your turn}
\begin{itemize}
\raggedright
  \item \emph{Say it:} \emph{gūzen}
  \item \emph{Say it:} \emph{gūzen}, once more
  \item \emph{Say it:} \emph{moshi gūzen kare ni attara, tegami o watashite kudasai} — If you happen to meet him, please give him the letter.
  \item \emph{Recall:} say the Japanese for probability, then the Japanese for to hit (the mark), then say \emph{gūzen} again
\end{itemize}

\begin{tcolorbox}[breakable,colback=teal!4,colframe=teal!35!black,title={Before you move on}]
What does \ja{ぐうぜん} mean? (\textquotedblleft{}By chance; a coincidence\textquotedblright{}.) Say it once more.
\end{tcolorbox}
\section[\texorpdfstring{\ja{きっかけ} (kikkake)}{きっかけ (kikkake)}]{\ja{きっかけ} (kikkake) — a chance; a start}

\label{lesson:JA-C510-kikkake}



\begin{quote}
Before the new one: say the Japanese for to hit (the mark), then the Japanese for by chance.
\end{quote}

\subsection*{You'll want to know: \ja{きっかけ}}
\textbf{\ja{きっかけ}} — \emph{kikkake} — \textquotedblleft{}a chance; a start\textquotedblright{}.

\begin{quote}
\textbf{\ja{もし} \ja{きっかけが} \ja{あったら、かれと} \ja{はなします。}} — \emph{moshi kikkake ga attara, kare to hanashimasu} — \textquotedblleft{}If I get the chance, I will talk to him.\textquotedblright{}
\end{quote}

\begin{cousinweb}
\textbf{\ja{きっかけ}} (\emph{kikkake}) is a native Japanese word (wago), inherited rather than borrowed from Chinese, and usually written in kana; its older origin is not settled (JMdict).
\end{cousinweb}

\begin{grammarlens}[title={What you've built}]
One more word for saying what happens if.
\end{grammarlens}

\subsection*{Your turn}
\begin{itemize}
\raggedright
  \item \emph{Say it:} \emph{kikkake}
  \item \emph{Say it:} \emph{kikkake}, once more
  \item \emph{Say it:} \emph{moshi kikkake ga attara, kare to hanashimasu} — If I get the chance, I will talk to him.
  \item \emph{Recall:} say the Japanese for to hit (the mark), then the Japanese for by chance, then say \emph{kikkake} again
\end{itemize}

\begin{tcolorbox}[breakable,colback=teal!4,colframe=teal!35!black,title={Before you move on}]
What does \ja{きっかけ} mean? (\textquotedblleft{}A chance; a start\textquotedblright{}.) Say it once more.
\end{tcolorbox}
\section[\texorpdfstring{\ja{ふと} (futo)}{ふと (futo)}]{\ja{ふと} (futo) — out of the blue; suddenly}

\label{lesson:JA-C510-futo}



\begin{quote}
Before the new one: say the Japanese for by chance, then the Japanese for a chance.
\end{quote}

\subsection*{You'll want to know: \ja{ふと}}
\textbf{\ja{ふと}} — \emph{futo} — \textquotedblleft{}out of the blue; suddenly\textquotedblright{}.

\begin{quote}
\textbf{\ja{もし} \ja{ふと} \ja{めが} \ja{さめたら、みずを} \ja{のみます。}} — \emph{moshi futo me ga sametara, mizu o nomimasu} — \textquotedblleft{}If I happen to wake up, I will drink some water.\textquotedblright{}
\end{quote}

\begin{cousinweb}
\textbf{\ja{ふと}} (\emph{futo}) is usually counted a native Japanese word (wago), inherited rather than borrowed from Chinese: JMdict marks it as usually written in kana and its rare kanji spellings as ateji, chosen for their sound, and no dictionary we use records an older form, so its history is not settled (JMdict).
\end{cousinweb}

\begin{grammarlens}[title={What you've built}]
One more word for saying what happens if.
\end{grammarlens}

\subsection*{Your turn}
\begin{itemize}
\raggedright
  \item \emph{Say it:} \emph{futo}
  \item \emph{Say it:} \emph{futo}, once more
  \item \emph{Say it:} \emph{moshi futo me ga sametara, mizu o nomimasu} — If I happen to wake up, I will drink some water.
  \item \emph{Recall:} say the Japanese for by chance, then the Japanese for a chance, then say \emph{futo} again
\end{itemize}

\begin{tcolorbox}[breakable,colback=teal!4,colframe=teal!35!black,title={Before you move on}]
What does \ja{ふと} mean? (\textquotedblleft{}Out of the blue; suddenly\textquotedblright{}.) Say it once more.
\end{tcolorbox}
\section[\texorpdfstring{\ja{いかす} (ikasu)}{いかす (ikasu)}]{\ja{いかす} (ikasu) — to make good use of}

\label{lesson:JA-C510-ikasu}



\begin{quote}
Before the new one: say the Japanese for a chance, then the Japanese for out of the blue.
\end{quote}

\subsection*{You'll want to know: \ja{いかす}}
\textbf{\ja{いかす}} — \emph{ikasu} — \textquotedblleft{}to make good use of\textquotedblright{}.

\begin{quote}
\textbf{\ja{もし} \ja{ちえを} \ja{いかしたら、いい} \ja{しごとが} \ja{みつかるでしょう。}} — \emph{moshi chie o ikashitara, ii shigoto ga mitsukaru deshō} — \textquotedblleft{}If you put that know-how to good use, you will probably find a good job.\textquotedblright{}
\end{quote}

Polite \textbf{\ja{いかします}}, \emph{ikashimasu}; a godan verb. Past (polite): \textbf{\ja{かれが} \ja{ちえを} \ja{いかしました。}}, \emph{kare ga chie o ikashimashita}.

\begin{cousinweb}
\textbf{\ja{いかす}} (\emph{ikasu}) is a native Japanese word (wago), inherited rather than borrowed from Chinese; its kanji was attached later as a kun'yomi, and its older origin is not settled (JMdict, KANJIDIC2).
\end{cousinweb}

\begin{grammarlens}[title={What you've built}]
One more word for saying what happens if.
\end{grammarlens}

\subsection*{Your turn}
\begin{itemize}
\raggedright
  \item \emph{Say it:} \emph{ikasu}
  \item \emph{Say it:} \emph{ikasu}, once more
  \item \emph{Say it:} \emph{moshi chie o ikashitara, ii shigoto ga mitsukaru deshō} — If you put that know-how to good use, you will probably find a good job.
  \item \emph{Recall:} say the Japanese for a chance, then the Japanese for out of the blue, then say \emph{ikasu} again
\end{itemize}

\begin{tcolorbox}[breakable,colback=teal!4,colframe=teal!35!black,title={Before you move on}]
What does \ja{いかす} mean? (\textquotedblleft{}To make good use of\textquotedblright{}.) Say it once more.
\end{tcolorbox}
\section[\texorpdfstring{\ja{めぐまれる} (megumareru)}{めぐまれる (megumareru)}]{\ja{めぐまれる} (megumareru) — to be blessed (with)}

\label{lesson:JA-C510-megumareru}



\begin{quote}
Before the new one: say the Japanese for out of the blue, then the Japanese for to make good use of.
\end{quote}

\subsection*{You'll want to know: \ja{めぐまれる}}
\textbf{\ja{めぐまれる}} — \emph{megumareru} — \textquotedblleft{}to be blessed (with)\textquotedblright{}.

\begin{quote}
\textbf{\ja{もし} \ja{てんきに} \ja{めぐまれたら、やまに} \ja{のぼります。}} — \emph{moshi tenki ni megumaretara, yama ni noborimasu} — \textquotedblleft{}If we are blessed with good weather, we will climb the mountain.\textquotedblright{}
\end{quote}

Polite \textbf{\ja{めぐまれます}}, \emph{megumaremasu}; an ichidan verb. Past (polite): \textbf{\ja{かれが} \ja{こどもに} \ja{めぐまれました。}}, \emph{kare ga kodomo ni megumaremashita}.

\begin{cousinweb}
\textbf{\ja{めぐまれる}} (\emph{megumareru}) is a native Japanese word (wago), inherited rather than borrowed from Chinese; its kanji was attached later as a kun'yomi, and its older origin is not settled (JMdict, KANJIDIC2).
\end{cousinweb}

\begin{grammarlens}[title={What you've built}]
One more word for saying what happens if.
\end{grammarlens}

\subsection*{Your turn}
\begin{itemize}
\raggedright
  \item \emph{Say it:} \emph{megumareru}
  \item \emph{Say it:} \emph{megumareru}, once more
  \item \emph{Say it:} \emph{moshi tenki ni megumaretara, yama ni noborimasu} — If we are blessed with good weather, we will climb the mountain.
  \item \emph{Recall:} say the Japanese for out of the blue, then the Japanese for to make good use of, then say \emph{megumareru} again
\end{itemize}

\begin{tcolorbox}[breakable,colback=teal!4,colframe=teal!35!black,title={Before you move on}]
What does \ja{めぐまれる} mean? (\textquotedblleft{}To be blessed (with)\textquotedblright{}.) Say it once more.
\end{tcolorbox}
